\documentclass[12pt]{article}

\usepackage[american]{babel}
\usepackage[margin=1in]{geometry} 
\usepackage{amsmath,amsthm,amssymb}
\usepackage{multicol}
\setlength{\columnsep}{1cm}
\usepackage{mathrsfs}
\usepackage{tikz}
\usepackage{color}
\usepackage{lipsum} % for filler text
\usepackage{fancyhdr}
\usepackage{graphicx}
\pagestyle{fancy}
\usepackage[svgnames]{xcolor}
\usepackage[colorlinks=true, allcolors=DeepPink]{hyperref}
\usepackage{hyperref}
\usepackage{float}
\usepackage{subcaption}
\usepackage{xcolor}
\usepackage{wrapfig}
\usepackage[backend=bibtex, style=alphabetic,]{biblatex}
\addbibresource{references.bib}

\newcommand{\N}{\mathbb{N}}
\newcommand{\Z}{\mathbb{Z}}
\newcommand{\ep}{\varepsilon}
\newcommand{\pa}{\partial}

\usepackage{titling}
\usepackage{hyperref}
\usepackage{longtable}

\setlength{\droptitle}{-8em}   % This is your set screw
\setlength{\headheight}{16pt}

\begin{document}

\title{AMATH 732 Group Project Research Notes}
\author{tosh AKA sarah}
\maketitle

\tableofcontents
\newpage

\section{Squid Hunting}

\subsection{Predatory behaviour of squid (Illex illecebrosus) feeding on surface swarms of euphausiids: \cite{nicol}}
\href{https://doi.org/10.1139/z85-003}{Link} \\
\begin{itemize}
    \item northern shortfin squid filmed feeding on surface swarms of krill
    \item captured prey via attacking with tentacles (its 1985)
    \item some nice diagrams of how they eat
    \item they attack in succession BUT there's no coordination between the squid, spatially nor temporally
    \item "the number of attacks per squid is thus independent of the number of squid in the vicinity of the swam, which suggests that the squid were neither hindered nor facilitated in their feeding by the close proximity of their conspecifics"
    \item different from other documented account of how they do this (head-first attack where tentacles not used in initial food capture)
    \item there was no indication that they favour krill on the periphery of the swarm instead of diving into the middle
\end{itemize}
VERDICT: pretty good ngl i thought this was all common knowledge but maybe common knowledge from 1985

\subsection{Predation behaviour of the European squid Loligo vulgaris: \cite{castilla}}
\href{https://doi.org/10.1007/s10164-020-00652-4}{Link}
\begin{itemize}
    \item european squid hunting strategy and body patterning when offered prawn and fish
    \item fish:
    \begin{itemize}
        \item tentacular lunge attack
        \item kill fish before eat
        \item always eaten in head-caudal fin direction
    \end{itemize}
    \item prawn:
    \begin{itemize}
        \item arm opening attack when it's smaller and less mobile (prawn)
        \item both arms and tentacles when prawn is further away
        \item eat alive
    \end{itemize}
    \item larger squid ate faster and consumed entire prey (better skills and bigger beak - older and more experienced)
    \item smaller squid rejected certain parts of fish prey
    \item really detailed description of how they eat in the results, and nice diagrams of body patterns
\end{itemize}
VERDICT: maybe too much detail for our purposes

\subsection{Observations of wild hunting behaviour and bioluminescence of a large deep-sea, eight-armed squid, Taningia danae: \cite{kubodera}}
\href{https://doi.org/10.1098/rspb.2006.0236}{Link}
\begin{itemize}
    \item dana octopus squid (no tentacles!)
    \item originally thought to be sluggish and neutrally buoyant, but it is not yay
    \item can change speeds and directions quickly using big fins
    \item short bright flashes from arm tip photophores (blind prey? measure distance?)
\end{itemize}
VERDICT: overall perhaps not relevant to our model since we are looking at a group of squid

\subsection{Cephalopods as Predators: A Short Journey among Behavioral Flexibilities, Adaptions, and Feeding Habits: \cite{villanueva}}
\href{https://doi.org/10.3389/fphys.2017.00598}{Link}
\begin{itemize}
    \item kind of like a lit review for how cephalopods eat omg?
\end{itemize}
VERDICT: could be worth to read in depth later

\subsection{Experiments of whether schooling by their prey affects the hunting behaviour of cephalopods and fish predators: \cite{neill}}
\href{https://doi.org/10.1111/j.1469-7998.1974.tb04385.x}{Link}
\begin{itemize}
    \item how much does the size of a fish school change the squid (and fish) hunting success?
    \item increasing the size of a school DID decrease hunting success rate
\end{itemize}
VERDICT: could be worth to read in depth later? idk how accurate it'll be to real squid behaviour

\subsection{Market squid (Doryteuthis opalescens) feeding habits off California during a period of high abundance: \cite{noord}}
\href{https://doi.org/10.1016/j.ecss.2026.109761}{Link}
\begin{itemize}
    \item market squid (in california) eat mostly krill, copepods, gastropods, teleosts (in that order) - first two are crustaceans, third is mollusc, fourth is fish
\end{itemize}
VERDICT: i mean good to know but not super relevant to our specific project (do krill school??)

\subsection{Predatory strategies of squid (Illex illecebrosus) attacking small and large fish: \cite{foyle}}
\href{https://doi.org/10.1080/10236248809378670}{Link}
\begin{itemize}
    \item how does fish size affect how the squid eats?
\end{itemize}
VERDICT: really good for my autism bc these are things i want to know BUT not relevant to our topic unfortunately :(

\subsection{A Kinematic Analysis of Tentacle Extension in the Squid Loligo Pealei: \cite{kier}}
\href{https://doi.org/10.1242/jeb.200.1.41}{Link}
\begin{itemize}
    \item how do squid use their tentacles to eat?
    \item only thing noteworthy is that the squid jets forward when attacking, which our model doesn't capture
    \item NO ACCESS??? :(
\end{itemize}
VERDICT: unfortunately not super relevant

\subsection{Food and feeding of jumbo squid Dosidicus gigas in the Gulf of California and adjacent waters after the 1997-98 El nino event: \cite{markaida}}
\href{https://doi.org/10.1016/j.fishres.2006.02.016}{Link}
\begin{itemize}
    \item humboldt squid mostly eat mesopelagic fish
    \item this was only for medium sized humboldt squid and after an el nino event so hard to say how reliable these are
\end{itemize}
VERDICT: thank goodness squid do actually eat fish

\subsection{Coordinated noctural behavior of foraging jumbo squid Dosidicus gigas: \cite{benoit-bird}}
\href{https://doi.org/10.3354/meps09664}{Link}
\begin{itemize}
    \item observations of humboldt squid eating at night
    \item they are very coordinated - constitute schooling!!!!
    \item "feeding events might define the anchor for group behaviour" and learning how they school could be important for understanding their biology
\end{itemize}
VERDICT: maybe worth a skim just qualitatively for background

\subsection{Feeding habits of juvenile Japanese common squid Todarodes pacificus: Relationship between dietary shift and allometric growth: \cite{uchikawa}}
\href{https://doi.org/10.1016/j.fishres.2013.07.001}{Link}
\begin{itemize}
    \item most important result: when these squid become over a certain size, their diet shifts from largely crustaceans to largely fish! so they really will just eat anything...
\end{itemize}
VERDICT: idk if worth reading

\subsection{Turning performance of brief squid Lolliguncula brevis during attacks on shrimp and fish: \cite{jastrebsky}}
\href{https://doi.org/10.1242/jeb.144261}{Link}
\begin{itemize}
    \item three behavioural phases: approach, strike, recoil
    \item different movements for shrimp and fish :O
    \item the turning is important wowww
\end{itemize}

\subsection{Faster is not always better: Turning performance trade-offs in the inshore squids Doryteuthis pealeii and Illex illecebrosus: \cite{ganley}}
\href{https://doi.org/10.1016/j.jembe.2023.151913}{Link}
\begin{itemize}
    \item how does fin size affect turning rate?
    \item kinematic analysis
    \item everything matters LOL is it better to have max angular velocity or minimum turn radius (you can only have one)
\end{itemize}


%%%%%%%%%%%%%%%%%%%%%%%%%%%%%%%%%%%%%%%%%%%%%%%%%%

\section{Fish Schooling}

\subsection{An individual based model of fish school reactions: predicting anitpredator behaviour as observed in nature: \cite{vabo}}
\href{https://doi.org/10.1046/j.1365-2419.1997.00037.x}{Link}
\begin{itemize}
    \item individual based model to predict behaviour of fish in general (case used norwegian herring)
    \item fish schools can "split, join, and vacuole" during a predator attack
    \item driving force of fish is based on attraction rules
    \item stochastic elements: herring do not have perfect info abt surroundings
    \item model also has isolation of individual fish during predator attack (maybe disrupting the school is important for predators?)
    \item model compared with observed predator-prey interactions from herring and orcas
    \item the fish have a Panic Distance - if there is a predator within the panic distance, they will go away from it
    \item predator perception length
    \item a bunch of parameters that we don't have in our model LOL
    \item speed variable
    \item results: ICs hardly affected
    \item results: a large variety of different fish school shapes observed
\end{itemize}
VERDICT: definitely worth a read

\subsection{Density distribution and size sorting in fish schools: an individual-based model: \cite{hemelrijk}}
\href{https://doi.org/10.1093/beheco/arh149}{Link}
\begin{itemize}
    \item what spatial patterns can be observed in fish schools as a result of different mechanisms? (ex. familiarity, body size, etc.)
    \item schools are usually oblong and density is highest near the front - potential explanation for why
    \item different school shapes
    \item individual fish know each other and stick with their friends?
    \item slightly different schooling mechanics: avoid fish close by, align with fish at intermediate distance, move towards fish further off (within range of vision)
    \item "size difference model"
    \item in nature they usually see large individuals in the centre of the school
\end{itemize}
VERDICT: definitely worth a read

%%%%%%%%%%%%%%%%%%%%%%%%%%%%%%%%%%%%%%%%%%%%%%%%%%%%

\section{Individual-Based Models}

\subsection{Individual-Based Modeling of Ecological and Evolutionary Processes: \cite{deangelis}}
\href{https://doi.org/10.1146/annurev.ecolsys.36.102003.152644}{Link}
\begin{itemize}
    \item How are IBMs useful and where are they used?
\end{itemize}
VERDICT: maybe worth a skim for some general info

\subsection{Growth Rate Variation and Larval Survival: Inferences from an Individual-Based Size Dependent Predation Model: \cite{rice}}
\href{https://doi.org/10.1139/f93-015}{Link}
\begin{itemize}
    \item Using and IBM to learn about baby fish mortality and how often they survive
    \item faster-growing fish survive for longer generally
\end{itemize}
VERDICT: not super relevant to our topic, our sample fish are already full grown

\section{Misc}

\subsection{Impact of cephalopods in the food chain and their interaction with the environment and fisheries: an overview: \cite{piatkowski}}
\href{https://doi.org/10.1016/S0165-7836(01)00226-0}{Link}
really good if we were studying climate change or doing the population model task

\subsection{Swarming and Behaviour in Antarctic Krill: \cite{tarling}}
\href{https://doi.org/10.1007/978-3-319-29279-3_8}{Link}
I skimmed this article to confirm that krill do in fact school! And they do :D So our model can be used with krill as well

\newpage
\printbibliography[heading=bibintoc]


\end{document}